% Ficha de estudio 2 del Taller virtual ATA EPE: dinámica del cuerpo rígido.
\documentclass[11pt,a4paper]{article}
% Preámbulo común de las fichas de estudio del Taller virtual ATA EPE.
\usepackage[spanish,es-noshorthands,es-tabla]{babel}
\usepackage[utf8]{inputenc}
\usepackage[T1]{fontenc}
\usepackage{lmodern}
\usepackage[margin=2.2cm]{geometry}
\usepackage[table]{xcolor}
\usepackage{amsmath,amssymb,bm}
\usepackage{tabularx,array,booktabs}
\usepackage{enumitem}
\usepackage{fontawesome5}
\usepackage[most]{tcolorbox}
\usepackage[hidelinks]{hyperref}
\usepackage{listings}

\definecolor{morado}{HTML}{2A1B7E}
\definecolor{amarillo}{HTML}{FFCF2E}
\definecolor{rojo}{HTML}{C4142C}
\definecolor{verde}{HTML}{1F9D55}
\definecolor{gris}{HTML}{5D5A73}

\newcolumntype{Y}{>{\raggedright\arraybackslash}X}
\renewcommand{\arraystretch}{1.25}
\setlist{itemsep=1pt,topsep=3pt}
\numberwithin{equation}{section}
\lstset{basicstyle=\ttfamily\small,columns=fullflexible,keepspaces=true,frame=single,rulecolor=\color{gris!40},
  literate={á}{{\'a}}1 {é}{{\'e}}1 {í}{{\'i}}1 {ó}{{\'o}}1 {ú}{{\'u}}1 {ñ}{{\~n}}1 {Ω}{{$\Omega$}}1
    {θ}{{$\theta$}}1 {λ}{{$\lambda$}}1 {Δ}{{$\Delta$}}1 {μ}{{$\mu$}}1 {·}{{$\cdot$}}1 {←}{{$\leftarrow$}}1 {§}{{\S}}1 {→}{{$\to$}}1 {≠}{{$\neq$}}1 {É}{{\'E}}1 {Á}{{\'A}}1 {∨}{{$\vee$}}1 {∧}{{$\wedge$}}1 {σ}{{$\sigma$}}1 {×}{{$\times$}}1 {≈}{{$\approx$}}1}

\newcommand{\rev}[1]{\widetilde{#1}}
\newcommand{\I}{\mathbf{I}}
\newcommand{\vect}[1]{\bm{#1}}
\newcommand{\A}{\mathcal{A}}

\newtcolorbox{meta}{colback=morado,colframe=morado,coltext=white,arc=3mm}
\newtcolorbox{lectura}{colback=morado!5,colframe=morado,fonttitle=\bfseries,title={\faBook~Lecturas}}
\newtcolorbox{ejercicio}[1]{colback=amarillo!10,colframe=amarillo!80!black,fonttitle=\bfseries,title={\faPencil*~Ejercicio #1}}
\newtcolorbox{control}{colback=rojo!5,colframe=rojo,fonttitle=\bfseries,title={\faQuestionCircle~Preguntas de control}}
\newtcolorbox{codigo}[1][De la ecuación al código]{colback=verde!6,colframe=verde,fonttitle=\bfseries,title={\faCode~#1}}
\newtcolorbox{cierre}{colback=verde!6,colframe=verde,fonttitle=\bfseries,title={\faFlagCheckered~Cierre del módulo}}
\newtcolorbox{nota}[1]{colback=amarillo!12,colframe=amarillo!80!black,fonttitle=\bfseries,title={\faLightbulb~#1}}

% \cabecera{número}{título}{subtítulo}
\newcommand{\cabecera}[3]{%
  \title{\textcolor{morado}{\textbf{Ficha #1. #2}}\\[3pt]\large Taller virtual ATA EPE · #3}
  \author{}\date{}}

\cabecera{2}{Masa, inercia y movimiento libre}{Módulo 2: dinámica del cuerpo rígido}

\begin{document}
\maketitle
\vspace{-2.5em}

\begin{meta}
\textbf{\textcolor{amarillo}{Meta del módulo.}} Calcular masa, centro de masa e inercia de cada pieza a partir de su
malla; escribir y resolver las ecuaciones de Newton–Euler en PGA; avanzar en el tiempo sin salirse del grupo de
movimientos rígidos, y entender qué integrador conserva qué. Cara pública: \textbf{laboratorio de giro} — lanzar una
pieza y ver su movimiento libre, incluido el efecto Dzhanibekov.
\end{meta}

\begin{lectura}
\begin{enumerate}[leftmargin=*]
  \item \textbf{Goldstein, Poole y Safko, \emph{Classical Mechanics}, cap.~5} (§5.1--5.7): tensor de inercia, ejes
        principales, ecuaciones de Euler, cuerpo libre.
  \item \textbf{Dorst y De Keninck, \emph{May the Forque Be with You}}, secciones de dinámica: la aplicación de
        inercia, el forque y la ecuación de movimiento en el marco del cuerpo.
  \item \textbf{Hairer, Lubich y Wanner, \emph{Geometric Numerical Integration}}, caps.~I y VI (§VI.1--VI.3):
        métodos simplécticos, por qué conservan la energía a largo plazo; y §IV.8 (métodos en grupos de Lie).
  \item Landau y Lifshitz, \emph{Mecánica}, §32--37: el trompo, Poinsot y la estabilidad de los ejes.
  \item Gunn, tesis doctoral, cap.~9: la misma dinámica en PGA, con más rigor geométrico.
  \item Mirtich, «Fast and accurate computation of polyhedral mass properties», \emph{Journal of Graphics Tools}
        1(2), 1996: masa e inercia de mallas.
\end{enumerate}
\end{lectura}

% =================================================================================================
\section{Propiedades de masa de una malla}

Una pieza es un sólido cerrado, cuya frontera es una malla de triángulos orientados hacia afuera. Cada triángulo
$(\vect a,\vect b,\vect c)$ forma con el origen un tetraedro con volumen con signo y matriz de covarianza
\begin{equation}
  V_t = \tfrac16\,\vect a\cdot(\vect b\times\vect c),\qquad
  \mathbf{C}_t = \det(\mathbf{A})\,\mathbf{A}\,\mathbf{C}_0\,\mathbf{A}^{\mathsf T},\qquad
  \mathbf{C}_0 = \tfrac{1}{120}\begin{pmatrix}2&1&1\\1&2&1\\1&1&2\end{pmatrix},
\end{equation}
con $\mathbf A = [\vect a\ \vect b\ \vect c]$ (columnas). Por el teorema de la divergencia, las partes que quedan
fuera del sólido se cancelan, y sumando sobre todos los triángulos:
\begin{equation}
  m = \rho\sum_t V_t,\qquad \vect c = \frac{\rho}{m}\sum_t V_t\,\frac{\vect a+\vect b+\vect c}{4},\qquad
  \mathbf C = \rho\sum_t \mathbf C_t - m\,\vect c\,\vect c^{\mathsf T},\qquad \mathbf I = \operatorname{tr}(\mathbf C)\,\mathbb 1 - \mathbf C.
\end{equation}
La última relación pasa de la covarianza (segundos momentos) al tensor de inercia respecto al centro de masa.

\begin{ejercicio}{2.1}
Aplica (1.1)--(1.2) a un cubo de lado 1 (12 triángulos) y verifica $I = m/6$ en los tres ejes. Luego a un cilindro
aproximado con 24 lados: ¿cuánto se aleja de $\tfrac12 mR^2$?
\end{ejercicio}

\subsection{Ejes principales}
$\mathbf I$ es simétrica: tiene una base ortonormal de vectores propios (ejes principales) con momentos
$I_1\le I_2\le I_3$. Se diagonaliza con el método de Jacobi (rotaciones de Givens que anulan un elemento fuera de la
diagonal por vez). En el marco principal, la inercia es diagonal y las ecuaciones se simplifican; el motor de la
pieza absorbe la rotación entre su marco de modelado y el principal.

% =================================================================================================
\section{Newton–Euler y su forma en PGA}

En el marco del cuerpo, con origen en el centro de masa y ejes principales:
\begin{align}
  m\,(\dot{\vect v} + \vect\omega\times\vect v) &= \vect f,\\
  \mathbf I\,\dot{\vect\omega} + \vect\omega\times\mathbf I\vect\omega &= \vect\tau .
\end{align}
En PGA, velocidad y fuerza son bivectores ($\Omega$, $F$) y la inercia es una aplicación lineal $\A$ de velocidades en
momentos. Las dos ecuaciones son una sola:
\begin{equation}
  \dot\Omega = \A^{-1}\bigl[\,F - \Omega\times\A[\Omega]\,\bigr].
\end{equation}

\begin{ejercicio}{2.2}
Con la definición de $\A$ de Dorst y De Keninck, escribe $\Omega = \omega_1e_{23}+\omega_2e_{31}+\omega_3e_{12}
+ v_1e_{01}+v_2e_{02}+v_3e_{03}$ (ajusta los signos a la convención de la ficha 0) y muestra que (2.3) contiene
exactamente a (2.1) y (2.2). ¿Dónde aparece $\vect\omega\times\vect v$?
\end{ejercicio}

\subsection{Fuerzas como rectas}
Una fuerza $\vect f$ aplicada en el punto $P$ es la recta que pasa por $P$ con dirección $\vect f$, escalada por
$\lvert\vect f\rvert$. Varias fuerzas se suman como bivectores; el resultado, en general, no es una recta sino un
tornillo de fuerza (\emph{wrench}): una fuerza más un par a lo largo de ella. La gravedad sobre una pieza es la recta
vertical por su centro de masa, de magnitud $mg$. Para usarla en (2.3), se lleva al marco del cuerpo:
$F_c = \rev M\,F_e\,M$.

\begin{ejercicio}{2.3}
Escribe la fuerza de gravedad sobre una pieza con centro de masa en $(x, y, z)$ como bivector. Muestra que su parte
euclídea es la fuerza y su parte ideal, el par respecto al origen.
\end{ejercicio}

% =================================================================================================
\section{El cuerpo libre}

Sin fuerzas, se conservan la energía $T = \tfrac12\,\Omega\vee\A[\Omega]$ y el momento (la recta $\A[\Omega]$ llevada
al espacio). El movimiento de rotación se ve en la construcción de Poinsot: el vector $\vect L$ fijo, el elipsoide de
energía rodando sobre un plano.

\subsection{El teorema del eje intermedio}
Linealizando (2.2) alrededor de un giro en el eje $k$ se obtiene, para las perturbaciones en los otros dos ejes $i,j$,
\begin{equation}
  \ddot\varepsilon = \omega_k^2\,\frac{(I_k - I_i)(I_j - I_k)}{I_iI_j}\,\varepsilon .
\end{equation}
Si $I_k$ es el intermedio, el factor es positivo y la perturbación crece: el giro es inestable. Es el «efecto
Dzhanibekov», y una viga del kit lo muestra (tiene tres momentos distintos).

\begin{ejercicio}{2.4}
Deduce (3.1) desde las ecuaciones de Euler. Calcula los momentos principales de la viga de 5 y estima cuánto tarda en
voltearse si se lanza girando a 3 vueltas por segundo con una perturbación de $1^\circ$.
\end{ejercicio}

% =================================================================================================
\section{Discretización: avanzar sin salir del grupo}

\subsection{El oscilador como banco de pruebas}
Para $\ddot x = -\omega^2x$ con paso $h$, cada método multiplica el estado $(x, v)$ por una matriz de amplificación:
\begin{align}
  \text{Euler explícito:}&\quad \begin{pmatrix}1 & h\\ -\omega^2h & 1\end{pmatrix},\quad \det = 1 + \omega^2h^2 > 1
    \ \text{(la energía crece)},\\
  \text{Euler semiimplícito:}&\quad \begin{pmatrix}1-\omega^2h^2 & h\\ -\omega^2h & 1\end{pmatrix},\quad \det = 1
    \ \text{(simpléctico; estable si } \omega h < 2).
\end{align}
El semiimplícito (primero la velocidad, luego la posición con la velocidad nueva) es simpléctico: no conserva la
energía exacta, sino una «energía modificada» cercana, y por eso no deriva en tiempos largos.

\begin{ejercicio}{2.5}
Verifica (4.1) y (4.2). Calcula los valores propios de (4.2) y muestra que tienen módulo 1 si $\omega h < 2$. ¿Qué pasa
con RK4: es estable?, ¿conserva la energía?
\end{ejercicio}

\subsection{El paso en el grupo}
La posición no es un vector sino un motor. El paso semiimplícito queda:
\begin{align}
  \Omega_{k+1} &= \Omega_k + h\,\A^{-1}\bigl[F_k - \Omega_k\times\A[\Omega_k]\bigr],\\
  M_{k+1} &= M_k\,\exp\!\bigl(-\tfrac h2\,\Omega_{k+1}\bigr).
\end{align}
Como la exponencial de un bivector es siempre un motor, $M$ nunca deja de ser un movimiento rígido: no hace falta
reortogonalizar matrices. Solo se renormaliza de vez en cuando por redondeo (ficha 0, §5.3).

\begin{nota}{El término giroscópico}
En (4.3), $\Omega\times\A[\Omega]$ se evalúa en el paso anterior (explícito). Para piezas largas que giran rápido,
eso inyecta energía. Si aparece en las pruebas, se resuelve ese término de forma implícita con una iteración de
Newton sobre la ecuación de Euler: es una mejora del módulo, no un cambio de método.
\end{nota}

\begin{codigo}[Qué programar]
\begin{tabularx}{\linewidth}{@{}l Y@{}}
\texttt{nucleo/masa.js} & \texttt{propiedadesDeMasa(vértices, triángulos, densidad)}: $m$, $\vect c$, $\mathbf I$
  (1.1--1.2); \texttt{ejesPrincipales(I)} por Jacobi.\\
\texttt{nucleo/cuerpo.js} & Clase \texttt{Cuerpo}: motor, velocidad en el cuerpo, masa, inercia principal, forques
  acumulados; \texttt{inercia}($\Omega$), \texttt{inerciaInversa(P)}, \texttt{energia()}, \texttt{momentoEnEspacio()}.\\
\texttt{nucleo/fuerzas.js} & \texttt{fuerzaEn(punto, vector)}, \texttt{gravedad(cuerpo)}.\\
\texttt{nucleo/integrador.js} & \texttt{paso(cuerpos, h)} con (4.3)--(4.4); subpasos.\\
\texttt{app/} & Laboratorio de giro: lanzar una pieza con la mano, con o sin gravedad.\\
\end{tabularx}
\end{codigo}

\begin{control}
\begin{enumerate}[leftmargin=*]
  \item ¿Por qué la suma de tetraedros con el origen da la masa correcta aunque el origen esté fuera de la pieza?
  \item ¿Por qué conviene trabajar en el marco principal del cuerpo?
  \item ¿Qué representa el término $\Omega\times\A[\Omega]$? ¿Por qué no aparece en el marco del espacio?
  \item ¿Por qué Euler explícito hace crecer la energía de un oscilador?
  \item ¿Qué significa que un método sea simpléctico y por qué importa en simulaciones largas?
  \item ¿Por qué la viga voltea al girar en su eje intermedio?
\end{enumerate}
\end{control}

\begin{cierre}
Ejercicios 2.1--2.5 · pruebas P1 (caída libre) y P11 (cuerpo libre: momento conservado con deriva menor que 0{,}1\,\%
en 10\,s) · masa e inercia de todas las piezas del catálogo · cara pública: el laboratorio de giro.
\end{cierre}

\end{document}
